\SourcePage{65}{70}

\section{Wave equation for a spin-1 particle}

In its rest frame, a spin-1 particle is described by a three-component wave function, namely a three-dimensional vector (such a particle is often called a \emph{vector particle}). As four-dimensional quantities, these components may originate either as the three spatial components of a spacelike 4-vector $\psi^\mu$, or as the mixed components of an antisymmetric rank-two 4-tensor $\psi^{\mu\nu}$. In the rest frame, the temporal component ($\psi^0$) of the former and the spatial components ($\psi^{ik}$) of the latter vanish\footnote{Anticipating a later result, the combined 4-vector $\psi_\mu$ and 4-tensor $\psi^{\nu\mu}$ correspond to the set of four-dimensional rank-two spinors $\xi^{\alpha\beta}$, $\eta_{\dot\alpha\dot\beta}$, and $\zeta^{\dot\alpha\dot\beta}$. Here $\xi^{\alpha\beta}$ and $\eta_{\dot\alpha\dot\beta}$ are symmetric spinors transformed into one another by inversion (see \S~19).}.

The wave equation, a differential relation between $\psi^\mu$ and $\psi^{\mu\nu}$, is specified by the relations
\begin{align}
i\psi_{\mu\nu}&=\widehat p_\mu\psi_\nu-\widehat p_\nu\psi_\mu,
\tag{14.1}\\
im^2\psi_\mu&=\widehat p^\nu\psi_{\mu\nu},
\tag{14.2}
\end{align}

\SourcePage{66}{71}
where $\widehat p=i\partial$ (A.~Proca, 1936). Apply $\widehat p^\mu$ to both sides of (14.2). The antisymmetry of $\psi_{\mu\nu}$ then gives
\begin{equation}
\widehat p^\mu\psi_\mu=0.
\tag{14.3}
\end{equation}

The quantity $\psi_{\mu\nu}$ can be eliminated from (14.1), (14.2) by substituting the first equation into the second. With (14.3), the result is
\begin{equation}
(\widehat p^2-m^2)\psi_\mu=0,
\tag{14.4}
\end{equation}
which again shows (compare \S~10) that $m$ is the particle mass. Thus a free spin-1 particle can be described using only one 4-vector $\psi^\mu$. Its components obey the second-order equation (14.4) together with the supplementary condition (14.3), which removes from $\psi^\mu$ the part belonging to spin~0.

In the rest frame, $\psi_\mu$ is independent of the spatial coordinates, and hence $\widehat p^0\psi_0=0$. At the same time $\widehat p^0\psi_0=m\psi_0$, so $\psi_0=0$ in that frame, as required. The components $\psi_{ik}$ vanish along with $\psi_0$.

A spin-1 particle may have either of two intrinsic parities, according as $\psi^\mu$ is a polar vector or a pseudovector. In the former case
\[
\widehat P\psi^\mu=(\psi^0,-\psi^i),
\]
whereas in the latter
\[
\widehat P\psi^\mu=(-\psi^0,\psi^i).
\]

Equations (14.1), (14.2) follow from a variational principle with the Lagrangian
\begin{align}
L={}&(1/2)\psi_{\mu\nu}\psi^{\mu\nu*}
-(1/2)\psi^{\mu\nu*}(\partial_\mu\psi_\nu-\partial_\nu\psi_\mu)
\notag\\
&-(1/2)\psi^{\mu\nu}(\partial_\mu\psi^*_\nu-\partial_\nu\psi^*_\mu)
+m^2\psi_\mu\psi^{\mu*}.
\tag{14.5}
\end{align}
Its independent generalized coordinates are $\psi_\mu$, $\psi^*_\mu$, $\psi_{\mu\nu}$, and $\psi^*_{\mu\nu}$\footnote{Had the variation been performed only with respect to $\psi_\mu$, with $\psi_{\mu\nu}$ assumed from the outset to be expressed through $\psi_\mu$ by (14.1), equation (14.3) would have had to be introduced as a supplementary condition unrelated to the variational principle.}.

For finding the energy--momentum tensor, formula (10.11) is not entirely convenient here: it would produce a nonsymmetric tensor requiring subsequent symmetrization. Instead one may use
\begin{equation}
\frac12 T_{\mu\nu}\sqrt{-g}
=-\frac{\partial}{\partial x^\lambda}
\frac{\partial\sqrt{-g}\,L}{\partial g^{\mu\nu}_{,\lambda}}
+\frac{\partial\sqrt{-g}\,L}{\partial g^{\mu\nu}},
\tag{14.6}
\end{equation}

\SourcePage{67}{72}
where $L$ is to be written in a form valid for arbitrary curvilinear coordinates (see II, \S~94). If $L$ contains only the components of the metric tensor $g_{\mu\nu}$ itself, and not their coordinate derivatives, the formula simplifies to
\[
T_{\mu\nu}=\frac{2}{\sqrt{-g}}\frac{\partial\sqrt{-g}\,L}{\partial g^{\mu\nu}}
=2\frac{\partial L}{\partial g^{\mu\nu}}-g_{\mu\nu}L
\]
(recall that $d\ln g=-g_{\mu\nu}dg^{\mu\nu}$).

The differentiation in (14.6) is not performed with respect to $\psi_\mu$ and $\psi_{\mu\nu}$. Consequently, these quantities need not be treated as independent when that formula is used: relation (14.1) can be applied at once, and (14.5) can be rewritten as
\begin{equation}
L=(-1/2)\psi_{\mu\nu}\psi^*_{\lambda\rho}g^{\mu\lambda}g^{\nu\rho}
+m^2\psi_\mu\psi^*_\nu g^{\mu\nu}.
\tag{14.7}
\end{equation}
It then follows that
\begin{align}
T_{\mu\nu}={}&-\psi_{\mu\lambda}\psi_\nu^{\ \lambda*}
-\psi^*_{\mu\lambda}\psi_\nu^{\ \lambda}
+m^2(\psi^*_\mu\psi_\nu+\psi^*_\nu\psi_\mu)
\notag\\
&+g_{\mu\nu}\left((1/2)\psi_{\lambda\rho}\psi^{\lambda\rho*}
-m^2\psi^*_\lambda\psi^\lambda\right).
\tag{14.8}
\end{align}

In particular, the energy density is the positive-definite expression
\begin{equation}
T_{00}=(1/2)\psi_{ik}\psi^*_{ik}+\psi_{0i}\psi^*_{0i}
+m^2(\psi_0\psi^*_0+\psi_i\psi^*_i).
\tag{14.9}
\end{equation}
The conserved 4-vector of current density is
\begin{equation}
j^\mu=i(\psi^{\mu\nu*}\psi_\nu-\psi^{\mu\nu}\psi^*_\nu).
\tag{14.10}
\end{equation}
It can be obtained from (12.12) by differentiating the Lagrangian (14.5) with respect to the derivatives $\partial_\mu\psi_\nu$. In particular,
\begin{equation}
j^0=i(\psi^{0k*}\psi_k-\psi^{0k}\psi^*_k)
\tag{14.11}
\end{equation}
and this quantity is not positive definite.

A plane wave normalized to one particle in the volume $V=1$ is
\begin{equation}
\psi_\mu=\frac{1}{\sqrt{2\varepsilon}}u_\mu e^{-ipx},
\qquad
u_\mu u^{\mu*}=-1,
\tag{14.12}
\end{equation}
where $u_\mu$ is a unit polarization 4-vector satisfying, by virtue of (14.3), the 4-transversality condition
\begin{equation}
u_\mu p^\mu=0.
\tag{14.13}
\end{equation}
Indeed, substitution of (14.12) in (14.9) and (14.11) gives
\[
T_{00}=-2\varepsilon^2\psi_\mu\psi^{\mu*}=\varepsilon,
\qquad
j^0=1.
\]

Unlike a photon, a massive vector particle has three independent polarization directions. Their corresponding amplitudes are given in (16.21).

\SourcePage{68}{73}
For partially polarized vector particles, the polarization density matrix is defined so that in a pure state it reduces to the product
\[
\rho_{\mu\nu}=u_\mu u^*_\nu
\]
(in analogy with (8.7) for photons). According to (14.12), (14.13), it obeys
\begin{equation}
p^\mu\rho_{\mu\nu}=0,
\qquad
\rho^\mu_{\ \mu}=-1.
\tag{14.14}
\end{equation}
For unpolarized particles, $\rho_{\mu\nu}$ must have the form $ag_{\mu\nu}+bp_\mu p_\nu$. Determining $a$ and $b$ from (14.14) yields
\begin{equation}
\rho_{\mu\nu}=-\frac13\left(g_{\mu\nu}-\frac{p_\mu p_\nu}{m^2}\right).
\tag{14.15}
\end{equation}

The vector-particle field is quantized in the same way as the scalar field, so there is no need to repeat the entire argument. The $\psi$-operators of the vector field are
\begin{align}
\widehat\psi_\mu
&=\sum_{\mathbf p\alpha}\frac{1}{\sqrt{2\varepsilon}}
\left(\widehat a_{\mathbf p\alpha}u_\mu^{(\alpha)}e^{-ipx}
+\widehat b^+_{\mathbf p\alpha}u_\mu^{(\alpha)*}e^{ipx}\right),
\notag\\
\widehat\psi^+_\mu
&=\sum_{\mathbf p\alpha}\frac{1}{\sqrt{2\varepsilon}}
\left(\widehat a^+_{\mathbf p\alpha}u_\mu^{(\alpha)*}e^{ipx}
+\widehat b_{\mathbf p\alpha}u_\mu^{(\alpha)}e^{-ipx}\right),
\tag{14.16}
\end{align}
where the index $\alpha$ labels the three independent polarizations.

The positive definiteness of (14.9) for $T_{00}$ and the indefiniteness of $j^0$ in (14.11) lead, just as in the scalar case, to Bose quantization.

There is a close connection between the properties of a genuinely neutral vector field and those of the electromagnetic field. A neutral vector field is described by a Hermitian $\psi$-operator:
\begin{equation}
\widehat\psi_\mu
=\sum_{\mathbf p\alpha}\frac{1}{\sqrt{2\varepsilon}}
\left(\widehat c_{\mathbf p\alpha}u_\mu^{(\alpha)}e^{-ipx}
+\widehat c^+_{\mathbf p\alpha}u_\mu^{(\alpha)*}e^{ipx}\right).
\tag{14.17}
\end{equation}
The Lagrangian of this field is
\begin{equation}
\widehat L
=\frac14\widehat\psi_{\mu\nu}\widehat\psi^{\mu\nu}
-\frac12\widehat\psi^{\mu\nu}
(\partial_\mu\widehat\psi_\nu-\partial_\nu\widehat\psi_\mu)
+\frac12m^2\widehat\psi_\mu\widehat\psi^\mu.
\tag{14.18}
\end{equation}

The electromagnetic field is obtained by taking $m=0$. The 4-vector $\psi^\mu$ then becomes the 4-potential $A^\mu$, while the 4-tensor $\psi^{\mu\nu}$ becomes the field tensor $F^{\mu\nu}$, related to the potential by definition (14.1). Equation (14.2) becomes $\partial^\nu\psi_{\mu\nu}=0$, which is the second pair of Maxwell equations. Condition (14.3) no longer follows from it, and is therefore no longer compulsory. Since the supplementary condition is absent, $\widehat\psi_\mu$ and $\widehat\psi_{\mu\nu}$ need not be regarded as independent ``coordinates'' in the Lagrangian,

\SourcePage{69}{74}
and (14.18) reduces to
\begin{equation}
\widehat L=-\frac14\widehat\psi_{\mu\nu}\widehat\psi^{\mu\nu}
\tag{14.19}
\end{equation}
in agreement with the familiar classical expression for the electromagnetic-field Lagrangian. This Lagrangian, together with the tensor $\widehat\psi_{\mu\nu}$, is invariant under an arbitrary gauge transformation of the ``potentials'' $\widehat\psi_\mu$. The connection of this fact with zero mass is evident: because of the term $m^2\widehat\psi_\mu\widehat\psi^\mu$, the Lagrangian (14.18) does not possess that invariance.
